% !TeX root = ../main.tex
\chapter{Role of Physics and Other Sciences}
\label{chap:physics-science}

To measure the currents of the various supplies, the load current is passed through a shunt resistor for each supply. Ohm's law, \(V=IR\), converts the sampled current into a voltage that can be measured by the LTC2945 PMIC. For example, the digital rail's 2 A current flowing
through a 39 m\(\Omega\) shunt produces a 78 mV differential voltage. The
resistor also converts part of the electrical energy into heat according to
Joule heating, \(P=I^2R\). That same 2 A operating point dissipates 0.156 W in
the shunt. These two equations create the main current-sensor design tradeoff:
a larger resistance produces more measurable voltage and better current
resolution, but also produces greater voltage drop and heat. The selected shunt
must therefore provide adequate ADC signal while remaining comfortably below its
power rating and while causing acceptably small voltage/IR drop to the monitored
rail.

Voltage measurement relies on the physics of electric potential difference and
on analog-to-digital (ADC) conversion. Each LTC2945 is referenced to its own local
power-supply domain via optocouplers, so it can measure the potential of that rail without
requiring all six rails to share a common ground. The converter represents an
analog input with a finite number of digital states. With the LTC2945's 25
mV-per-count resolution, quantization sets the bar for the lower measurable limit. The current channel similarly resolves 25 \(\mu\)V per count across the shunt. Mathematics is required not only to convert the ADC code into volts and amps but also to determine whether the selected shunt and expected signal range make use of the full
ADC span. The calculations in Appendix~\ref{app:design-calculations} explicitly
convert between current, shunt voltage, ADC counts, and measurement resolution.

Galvanic isolation is another important physical concept in the project. Two
circuits can exchange information without sharing a conductive DC path by using
an isolator. In the ACPL-064L, an electrical input drives an LED, which is then detected by a photodiode on the other side, forming an isolated barrier between the two sides of the circuit. The signal crosses the isolation barrier while the two circuit grounds remain
electrically independent. This is important because directly connecting the
returns of floating supplies can create ground loops or force current through a
path that the original power system did not intend. The mathematical treatment
of the isolation interface includes logic voltage limits, LED/interface current,
and the allowable voltage across the isolation barrier. Although the monitored
rails are only in the tens of volts, the isolation components have ratings in excess of 100V, providing ample headroom.

The digital communication buses also depend on basic circuit behavior. An
I\textsuperscript{2}C bus uses open-drain I/O, meaning that devices pull a
line low while a pull-up element or active current source returns it high. Bus
capacitance and pull-up strength determine the line's electrical rise time. Pull ups that are too weak, or too much bus capacitance can prevent a logic-high level from being
reached quickly enough at the selected clock rate. The TCA9803 buffer used in
the design includes an active current source on one side, which eliminates the need to be concerned with setting RC time constants so long as bus capacitance remains below the max for the driver.

Power conversion on the controller hardware also uses electromagnetic energy
storage. The TPS62912 is a switching buck regulator: it alternately connects and
disconnects the input through switching devices while an inductor and capacitors
store and release energy to produce a lower average DC output. The basic
relationships among voltage, current, inductance, capacitance, switching
frequency, ripple, and load determine output quality. For this project the regulator is required to produce a 3.3 V logic rail from a higher input while supplying the display, Ethernet, and interface electronics. The preliminary load estimate is compared with the regulator's output-current rating so the design has sufficient current and thermal margin rather than selecting the regulator only from its nominal output voltage.

Finally, the project relies on mathematics to turn these physical principles
into engineering decisions. Kirchhoff's and Ohm's laws are used to solve the
prototype test network; \(P=VI\), \(P=I^2R\), and related equations are used for
component heating and power budgeting; ADC scale factors are used to convert
between raw counts and engineering units; inequalities define acceptable alarm
regions; and tolerance analysis estimates how resistor error and converter error
propagate into the reported current. The firmware then implements these same
relationships numerically. In this sense the software is not separate from the
physics of the hardware: it is the mathematical model that converts measured
electrical quantities into the voltage, current, alarm, and status information
presented to the user and to EPICS.
